\section{Methodology}

%------------------------------------------------------------------
\subsection{Implementation Stack}
%------------------------------------------------------------------

Models are implemented in PyTorch~\cite{paszke2019pytorch}, with \texttt{torchaudio}~\cite{hwang2023torchaudio}
for audio processing and CTC beam-search decoding. The character vocabulary is provided by the
\texttt{facebook/wav2vec2-base} tokenizer via Hugging Face \texttt{transformers}~\cite{wolf-etal-2020-transformers}.
CER and WER are computed using \texttt{torchmetrics}~\cite{Detlefsen2022}.

%------------------------------------------------------------------
\subsection{Shared CNN Feature Extractor}
\label{sec:cnn}
%------------------------------------------------------------------

All model variants share a two-layer 2-D convolutional front-end
(\(\text{Conv2D} \to \text{BatchNorm2D} \to \text{Clipped ReLU}\), where
\(\sigma(x) = \min(\max(x,0),20)\)), clipping activations at 20 to prevent exploding
values that can destabilise CTC training.
The first layer uses a \(41\times11\) kernel with stride \(2\times2\) and padding \(20\times5\);
the second uses a \(21\times11\) kernel with stride \(2\times1\) and padding \(10\times5\)
(frequency\(\times\)time).

%------------------------------------------------------------------
\subsection{Temporal Encoder Variants}
%------------------------------------------------------------------

\subsubsection{GRU / LSTM}

Both the GRU and LSTM use gated recurrent units to mitigate vanishing gradients; the LSTM additionally maintains an explicit cell state for stronger long-range memory.
Both variants use \(L=3\) stacked layers with hidden size 512, LayerNorm between layers, and dropout \(p=0.3\).
We evaluate unidirectional and bidirectional GRU; the unidirectional variant is followed by an additional \emph{look-ahead convolution} (kernel size 41, right-padded) for bounded future context.

\subsubsection{Conformer}

The Conformer~\cite{gulati2020conformer} augments a multi-head self-attention (MHSA) transformer
with a depthwise convolutional module interleaved between each attention layer.
Self-attention captures global dependencies across the full sequence in parallel,
while the convolution module provides a local-pattern inductive bias that pure transformers lack.
The \(\mathcal{O}(n^2)\) memory cost of self-attention is acceptable for the utterance lengths
in LJ Speech.
We substitute relative positional encoding with absolute sinusoidal encoding, which suffices for this single-speaker corpus and reduces compute.
The shared CNN front-end serves as the convolutional subsampling stage.
We explored 2 Conformer configurations:

\begin{table}[h]
\label{tab:conformers}
  \centering
  \begin{tabular}{lll}
    \toprule
    Latent Dimension & \#Conformer Blocks & \#Parameters (M) \\
    \midrule
    256 & 12 & 18 \\
    144 & 10 & 5 \\
    \bottomrule
  \end{tabular}
\end{table}

and both with 4 self-attention heads, dropout $=0.1$, and depthwise convolution kernel size $=15$.

%------------------------------------------------------------------
\subsection{Training}
%------------------------------------------------------------------

All models are trained end-to-end with Adam + OneCycleLR, mixed-precision (using \texttt{torch.amp}),
and gradient clipping. The best checkpoint is selected by validation WER.

\subsubsection{Replication}

To obtain statistically reliable estimates of model performance and to avoid reporting results that
are artefacts of a particular random initialisation, each model configuration is trained under
three seeds (1508, 2603, and 9102). Before each run, \texttt{torch.manual\_seed}, \texttt{numpy.random.seed},
and Python's built-in \texttt{random.seed} are all set to the chosen seed, allowing our results to be reproducible.
Performance metrics (CER and WER) are reported as the mean \(\pm\) 1 standard deviation across the
three seeds.
This guards against over-stating the benefit of any single architecture choice and enables
fair pairwise comparison between model variants.